\theory{Wheel theory}{Can algebra make division total, including division by zero, without silently pretending ordinary inverse laws still hold?}{A wheel is an algebraic structure in which division is always defined. It extends a commutative ring or semiring with an inverse-like operation and a special value for $0/0$.}{Universal algebra, computer algebra, formal specification, and division by zero.}{Totalized division and changed algebraic identities provide a consistent setting in which division remains a defined operation at zero.}

\paragraph{Motivation and history}
In ordinary arithmetic, $a/b$ is undefined when $b=0$. Computer systems may need every operation to return a value. Jesper Carlstrom developed the theory of wheels in the early 2000s \cite{wheel_ref1}.

The name comes from viewing the real projective line as a circle and adding a bottom point representing $0/0$. A complex projective line is a sphere; adding the extra point gives a related geometric picture.

\end{historylist}
\section*{3. Theory}
\subsection*{Core theory}
\paragraph{Definition and changed rules}
A wheel has a set, constants $0$ and $1$, commutative associative addition and multiplication, and a unary operation written $/x$. Division is shorthand for $x/y=x\cdot(/y)$. The inverse operation is involutive:
\[
//x=x.
\]
It is multiplicative, and additional axioms modify familiar distributive rules so that expressions remain defined when a denominator is zero.

In a wheel, $x/x$ need not equal $1$, and $0x$ need not equal $0$ in the general case. These changes prevent illegal cancellation from turning $0/0$ into an ordinary number. The special value $0/0$ behaves like a bottom element \cite{wheel_ref2}.

\paragraph{Uses and interpretation}
Wheels provide a total algebraic data type for rational expressions. A computer algebra system can manipulate an expression without repeatedly checking whether a denominator vanishes; the result records the exceptional value instead of losing it. A wheel agrees with ordinary arithmetic where denominators are nonzero.

\paragraph{A worked example}
Consider the expression
\[
 \frac{x^2-1}{x-1}.
\]
For an ordinary number with $x\ne1$, factoring gives
\[
 \frac{(x-1)(x+1)}{x-1}=x+1.
\]
At $x=1$, the original expression is $0/0$, while the shortened expression
has value $2$. Ordinary algebra must therefore remember the condition
$x\ne1$; cancelling first and substituting later changes the question.

A wheel keeps the exceptional value visible. The original expression is
assigned the wheel's value for $0/0$, while the expression $x+1$ remains
equal to $2$ at $x=1$. The two expressions agree where cancellation is
legal, but they are not silently declared equal everywhere. This is the
practical benefit of totalized arithmetic: a symbolic program can continue
its calculation without confusing an exceptional value with an ordinary
number \cite{wheel_ref1,wheel_ref2}.

\section*{4. Important theorems and results}
\begin{enumerate}[leftmargin=*,itemsep=0.35em]
\item \textbf{Total division.} Every pair $x,y$ has a defined quotient $x/y$.

\item \textbf{Involution.} Applying the inverse operation twice returns the original element.

\item \textbf{Extension principle.} A commutative ring or semiring can be extended to a wheel under suitable constructions.

\item \textbf{Projective model.} The real projective line with an added bottom point gives the geometric picture behind a real wheel.

\item \textbf{Modified cancellation.} Identities such as $x/x=1$ are restricted because they fail at exceptional values.

The results describe how a wheel extends ordinary arithmetic while retaining a total division operation. Involution supplies the inverse-like structure, and the extension principle embeds familiar rings or semirings into the wider system. Modified cancellation is essential: division is always defined, but ordinary field identities cannot be applied at exceptional values without checking the wheel's rules.
\end{enumerate}
\noindent\emph{Source for these results:} \cite{wheel_ref1}.

\theoryapplications
\theoryconnections{wheel_theory}{%
\item Universal algebra asks which equations an operation must satisfy, while rings and semirings provide the ordinary addition and multiplication that a wheel extends.
\item The key change is to make division total: instead of leaving $1/0$ undefined, the structure assigns it an explicit value and records which familiar cancellation laws must be weakened.
\item Formal specification is needed because software must know exactly how exceptional values propagate \cite{wheel_ref1,wheel_ref2}.
}{%
\item Wheel theory helps computer algebra systems manipulate expressions without stopping whenever an intermediate denominator becomes zero.
\item A symbolic program can carry the exceptional value through addition, multiplication, and inversion, then report the result under the wheel's rules.
\item The connection is useful for formal verification because every expression has a defined denotation, but ordinary field identities may not be used unless their nonzero assumptions have been checked \cite{wheel_ref1,wheel_ref3}.
}
\section*{Glossary}
\begin{description}[leftmargin=*,style=nextline,itemsep=0.3em]
\item[\textbf{Wheel.}] An algebraic structure in which division is always defined, including division by zero, with modified cancellation laws.
\item[\textbf{Total division.}] The property that every pair of elements $x, y$ has a defined quotient $x/y$, even when $y = 0$.
\item[\textbf{Inverse operation.}] A unary operation written $/x$ that acts like a reciprocal; division is defined as $x/y = x \cdot (/y)$.
\item[\textbf{Involution.}] The property that applying the inverse operation twice returns the original element: $//x = x$.
\item[\textbf{Bottom element.}] The special value representing $0/0$ in a wheel, which behaves differently from ordinary numbers.
\item[\textbf{Extension principle.}] A construction that embeds a commutative ring or semiring into a wheel while preserving ordinary arithmetic where defined.
\item[\textbf{Modified cancellation.}] The restriction that identities like $x/x = 1$ cannot be applied at exceptional values without checking the wheel's rules.
\item[\textbf{Projective line.}] The real or complex projective line with an added bottom point, providing the geometric picture behind a wheel.
\item[\textbf{Commutative ring.}] An algebraic structure with addition and multiplication satisfying usual axioms; a wheel extends this to make division total.
\item[\textbf{Semiring.}] An algebraic structure like a ring but without requiring additive inverses; wheels can extend these.
\item[\textbf{Cancellation law.}] The rule that $ax=bx$ implies $a=b$; in wheels such laws must be weakened because division by zero is defined.
\end{description}
\section*{7. Sample Questions}
\emph{Exercise reference:} \cite{wheel_ref1}. The cited edition is the source for the exercise set; consult its Exercises section for the official statement and numbering.
\begin{enumerate}[leftmargin=*,itemsep=0.9em]
\item \textbf{Exercise 1.} In a wheel, the inverse operation satisfies $//x = x$ (involution). Using this, show that $/(1/a) = a$ for any element $a$. \emph{Solution.}$/(1/a) = //a = a$, applying involution twice. More carefully: the inverse operation is a unary operation $/(\cdot)$, so $/(1/a)$ means we apply $/$ to the element $1/a$. By involution, $/(1/a) = //a = a$.

\item \textbf{Exercise 2.}$\mathbb{R} \cup \{\infty, \bot\}$ forms a wheel (the extended real wheel). For $x, y \in \mathbb{R}$ with $y \neq 0$, compute $x/y$ in the wheel and confirm it agrees with ordinary division. \emph{Solution.}$x/y = x \cdot (/y) = x \cdot (1/y) = x/y$, which is the ordinary real quotient. When $y = 0$, $x/0 = x \cdot (/0)$ where $/0$ is the special element $\infty$ in the wheel. Thus $x/0$ is defined and equals $\infty$ (for $x \neq 0$), or $\bot$ (for $x = 0$, since $0/0 = \bot$). The wheel agrees with ordinary arithmetic where the denominator is nonzero.

\item \textbf{Exercise 3.} In a wheel, $x/x$ need not equal $1$ for all $x$. Show that in the real wheel $\mathbb{R} \cup \{\infty, \bot\}$, $0/0 = \bot \neq 1$. \emph{Solution.}$0/0 = 0 \cdot (/0) = 0 \cdot \infty$. In the real wheel, this is defined as $\bot$ (the bottom element), not $1$. This prevents the contradiction: if $0/0 = 1$, then from $0 \cdot a = 0 \cdot b$ we could conclude $a = b$ by canceling $0/0$, which is false. The wheel assigns $\bot$ to maintain consistency.

\item \textbf{Exercise 4.} Consider the expression $\frac{x^2 - 4}{x - 2}$. In ordinary algebra, this equals $x + 2$ for $x \neq 2$. In the real wheel, what is the value at $x = 2$? \emph{Solution.}$\frac{x^2 - 4}{x - 2} = \frac{(x-2)(x+2)}{x - 2}$. At $x = 2$: the numerator is $0$ and the denominator is $0$, so the value is $0/0 = \bot$. The wheel correctly distinguishes $\frac{x^2-4}{x-2}$ from $x + 2$ at $x = 2$: the former gives $\bot$, while the latter gives $4$.

\item \textbf{Exercise 5.}$\mathbb{CP}^1$ (the Riemann sphere) is the complex projective line. Explain how it relates to the complex wheel by adding a bottom element. \emph{Solution.}$\mathbb{CP}^1 = \mathbb{C} \cup \{\infty\}$ is the one-point compactification of $\mathbb{C}$, and forms a wheel with $\infty$ playing the role of $1/0$ and $\bot$ playing the role of $0/0$. The complex wheel extends $\mathbb{C}$ by adding two special elements: $\infty$ (for $z/0$ with $z \neq 0$) and $\bot$ (for $0/0$). The Riemann sphere $\mathbb{CP}^1$ captures the $\infty$ part; the full complex wheel additionally includes $\bot$ to make all divisions total.

\end{enumerate}
\section*{8. References}
\begin{thebibliography}{9}
\bibitem{wheel_ref1} J. Carlstr{\"o}m, \emph{Wheels: On Division by Zero}, Licentiate thesis, Chalmers University of Technology, 2001.
\bibitem{wheel_ref2} J. A. Bergstra and J. V. Tucker, \emph{The Rational Numbers as an Abstract Data Type}, J.\ ACM, 2007.
\bibitem{wheel_ref3} A. Setzer, \emph{Wheels}, Manuscript, University of Nottingham, 2001.
\end{thebibliography}
